\documentclass[12pt]{article}

\usepackage{sbc-template}
\usepackage{graphicx,url}
\usepackage[utf8]{inputenc}
\usepackage[brazilian]{babel}
\usepackage{minted}
\usepackage{float}
\usepackage{subcaption}
\usepackage{booktabs}
\usepackage{url}
\usepackage{pdflscape}
\usepackage{hyperref}
\usepackage{tikz}
\usetikzlibrary{positioning, arrows.meta, shapes.geometric}

\renewcommand{\listingscaption}{Listagem}

\NewCommandCopy{\svtt}{\texttt}
\renewcommand\texttt[1]{\svtt{\small#1}}

\setminted{
  style=friendly,        % tema de cores
  fontsize=\footnotesize,% tamanho adequado para colunas SBC
  % linenos=true,          % numeração de linhas
  breaklines=true,       % quebra linhas longas
  frame=lines,           % linhas acima e abaixo
  framesep=4mm
  % baselinestretch=1.1,
  % xleftmargin=2em,
  % numbersep=6pt,
  % numberfirstline=true
}
     
\sloppy

\title{Jass (Jaime's Home Assistant Language)}

\author{Jaime Antonio Daniel Filho\inst{1} }

\address{Curso de Ciência da Computação -- \\ Universidade Federal de Santa Maria (UFSM)
  \email{jafilho@inf.ufsm.br}
}

\begin{document} 

\maketitle

\section{Introdução}

O Home Assistant consolidou-se como uma das principais plataformas de automação residencial, permitindo integrar e controlar dispositivos de diferentes fabricantes em um único ambiente. Entretanto, a criação de automações é baseada em arquivos YAML (\textit{YAML Ain't Markup Language}), cuja estrutura pode se tornar extensa e complexa mesmo para tarefas simples. Uma automação básica, como acender uma luz ao abrir uma porta, exige o conhecimento de diversos elementos da sintaxe, como gatilhos, condições, ações, serviços e parâmetros específicos de cada dispositivo. Além disso, muitos erros só são percebidos durante a execução da automação, seja por um nome de serviço incorreto, um parâmetro incompatível ou um tipo de valor inválido, dificultando a identificação da causa do problema e tornando o processo de desenvolvimento menos acessível para usuários sem experiência técnica.

Com o objetivo de reduzir essa complexidade, este trabalho apresenta o \texttt{jass}, uma \emph{Domain-Specific Language} (DSL) capaz de transpilar descrições de alto nível para o formato YAML utilizado pelo Home Assistant. A proposta é oferecer uma sintaxe mais simples e intuitiva, abstraindo detalhes estruturais da plataforma e permitindo que o usuário descreva apenas o comportamento desejado da automação. Além disso, a linguagem incorpora mecanismos de validação léxica, sintática e semântica, possibilitando que erros, como construções inválidas ou operações incompatíveis com determinados dispositivos, sejam identificados ainda durante a compilação. Para isso, o \texttt{jass} define um conjunto reduzido de comandos cujo significado é interpretado de acordo com o tipo do dispositivo envolvido, cabendo ao compilador verificar a consistência dessas operações e gerar automaticamente a configuração YAML correspondente.


A Listagem~\ref{lst:jass_example1} e a Listagem~\ref{lst:jass_example1_yaml} ilustram, respectivamente, a automação completa em
\texttt{jass} (8 linhas) e o YAML equivalente que o compilador gera (22
linhas). Note como \texttt{turn\_on(light.varanda,\ ...)} se
expande para o serviço \texttt{light.turn\_on} com a estrutura
\texttt{target}/\texttt{data}, e como \texttt{delay(5m)} vira a string
\texttt{"00:05:00"} no formato exigido pela plataforma.

\begin{listing}[H]
    \centering
    \begin{minted}{swift}
automation "Luz da varanda" {
    when (binary_sensor.porta changes to "on")
    do {
        turn_on(light.varanda, brightness=180)
        delay(5m)
        turn_off(light.varanda)
    }
}
    \end{minted}
    
    \caption{Uma automação básica em \texttt{jass}.}
    \label{lst:jass_example1}
\end{listing}

\begin{listing}[H]
    \centering
    \begin{minted}{yaml}
- alias: "Luz da varanda"
  mode: "single"
  triggers:
    - trigger: state
      entity_id: binary_sensor.porta
      to: "on"
  actions:
    - action: light.turn_on
      target:
        entity_id: light.varanda
      data:
        brightness: 180
    - delay: "00:05:00"
    - action: light.turn_off
      target:
        entity_id: light.varanda
    \end{minted}
    
    \caption{A saída YAML da automação anterior.}
    \label{lst:jass_example1_yaml}
\end{listing}

O restante do relatório está organizado da seguinte forma: A Seção
3 descreve a arquitetura em \emph{pipeline}; a Seção 4 define a
linguagem de forma conceitual e a Seção 5 a formaliza como gramática
livre de contexto. As Seções 6 a 8 detalham as fases de análise e a Seção 9 consolida a estratégia de detecção e
relato de erros. A Seção 10 trata da geração de código, a Seção 11 reúne
exemplos completos e a Seção 12 conclui, sintetizando os resultados e abordando as limitações.

\section{Instruções para Execução}

O projeto usa o gerenciador \href{https://docs.astral.sh/uv/}{uv}. A partir do diretório do
trabalho, instale as dependências e o comando \texttt{jass} (Listagem~\ref{lst:uv_install}):

\begin{listing}[H]
    \centering
    \begin{minted}{bash}
uv sync
\end{minted}
    \caption{Comando para configurar o projeto.}
    \label{lst:uv_install}
\end{listing}

Para transpilar um arquivo \texttt{.jass} para YAML (Listagem~\ref{lst:uv_run}):

\begin{listing}[H]
    \centering
    \begin{minted}{bash}
uv run jass caminho/para/automacao.jass
\end{minted}
    \caption{Comando para transpilar um arquivo .jass para YAML.}
    \label{lst:uv_run}
\end{listing}

Para gravar o resultado em arquivo, ao invés de mostrar no terminal, com \texttt{-o} (Listagem~\ref{lst:uv_output}):

\begin{listing}[H]
    \centering
    \begin{minted}{bash}
uv run jass caminho/para/automacao.jass -o automacao.yaml
    \end{minted}
    \caption{Comando para gravar o resultado em arquivo.}
    \label{lst:uv_output}
\end{listing}

Para inspecionar uma fase intermediária do compilador, com
\texttt{-\/-emit} (Listagem~\ref{lst:uv_emit}), especificando \texttt{tokens}, \texttt{parse},
\texttt{ast}, \texttt{check} ou \texttt{yaml}:

\begin{listing}[H]
    \centering
    \begin{minted}{bash}
# imprime a AST tipada
uv run jass caminho/para/automacao.jass --emit ast   

# roda a análise semântica
uv run jass caminho/para/automacao.jass --emit check 
    \end{minted}
    \caption{Comando para inspecionar uma fase intermediária do compilador.}
    \label{lst:uv_emit}
\end{listing}

Em caso de erro, os diagnósticos são impressos no \texttt{stderr} e o
programa termina com código não nulo.

\section{Visão Geral da Arquitetura}

O \texttt{jass} é organizado como um compilador em \emph{pipeline}, ou
seja, em uma sequência de fases em que cada uma consome o artefato
produzido pela anterior e gera o artefato da seguinte. Essa estrutura
mantém as responsabilidades separadas e torna o sistema fácil de testar
e inspecionar fase a fase. A Figura~\ref{fig:fluxo_compilacao} resume o fluxo, e a Tabela~\ref{tab:fases_modulos}
associa cada fase ao módulo que a implementa e ao artefato que ela
produz.

\begin{figure}[htpb]
    \centering
    \begin{tikzpicture}[
        % Configuração de espaçamento e estilo das setas
        node distance = 1.2cm,
        auto,
        % Estilo padrão das caixas de texto
        box/.style = {
            rectangle, 
            draw=black, 
            thick,
            align=center, % Permite quebra de linha com \\
            minimum width=4.5cm, 
            minimum height=0.8cm, 
            rounded corners=3pt
        },
        % Estilo das setas
        arrow/.style = {
            -{Stealth[scale=1.2]}, 
            thick
        }
    ]

    % --- Definição dos Nodos (Vertical) ---
    \node [box] (src) {Texto-fonte\\(.jass)};
    \node [box, below=of src] (tok) {Tokens};
    \node [box, below=of tok] (tree) {Árvore de\\derivação};
    \node [box, below=of tree] (ast) {AST tipada};
    \node [box, below=of ast] (chk) {Diagnósticos +\\tabela de símbolos};
    \node [box, below=of chk] (model) {Modelo de saída\\(dicts/listas)};
    \node [box, below=of model] (yaml) {YAML do\\Home Assistant};

    % --- Definição das Arestas com Labels ---
    \draw [arrow] (src) -- node[right] {análise léxica} (tok);
    \draw [arrow] (tok) -- node[right] {análise sintática} (tree);
    \draw [arrow] (tree) -- node[right] {transformação} (ast);
    \draw [arrow] (ast) -- node[right] {análise semântica} (chk);
    \draw [arrow] (chk) -- node[right] {geração de código} (model);
    \draw [arrow] (model) -- node[right] {serialização} (yaml);

    \end{tikzpicture}
    
    % Legenda e Label padrão SBC (abaixo da figura)
    \caption{Fluxo do compilador.}
    \label{fig:fluxo_compilacao}
\end{figure}

\begin{table}[H]
\centering
\begin{tabular}{l l l}
\toprule
\textbf{Fase} & \textbf{Módulo}\\
\midrule
Análise léxica & \texttt{parser.py} (Lark) \\
Análise sintática & \texttt{parser.py} (Lark, LALR) \\
Transformação & \texttt{transform.py}, \texttt{ast.py}, \texttt{values.py} \\
Análise semântica & \texttt{analyzer.py}, \texttt{registry.py} \\
Geração de código & \texttt{codegen.py} \\
Orquestração / diagnósticos & \texttt{cli.py}, \texttt{errors.py}  \\
\bottomrule
\end{tabular}
\caption{Fases do compilador e módulos correspondentes.}
\label{tab:fases_modulos}
\end{table}


Os artefatos intermediários são, em ordem: o texto-fonte; a sequência de tokens, unidades léxicas como
palavras-chave, identificadores e literais; a árvore de derivação
(\emph{parse tree}), que reflete a estrutura gramatical da
entrada; a AST (\emph{Abstract Syntax Tree}), uma representação
tipada e enxuta, na qual cada nó é um objeto Python com tipo próprio e
posição de origem; o modelo de saída, uma árvore de dicionários e listas
Python já na forma esperada pelo Home Assistant; e, por fim, o YAML
serializado. A análise semântica é uma fase à parte: percorre a AST para
preencher a tabela de símbolos e validar a entrada, produzindo a decisão de que o programa está correto ou a
lista de erros.

\subsection{A ferramenta Lark}

As duas primeiras fases são construídas sobre o Lark, uma biblioteca
Python para construção de analisadores. O Lark recebe uma gramática em
notação EBNF e oferece, a partir dela, o \emph{lexer} (analisador
léxico) e o \emph{parser} (analisador sintático), além de utilitários
para o tratamento de erros. O \texttt{jass} o emprega na configuração
LALR(1) com \emph{lexer} contextual, aproveitando:

\begin{itemize}
\tightlist
\item
  a propagação automática de linha e coluna para cada nó, base dos
  diagnósticos localizados;
\item
  as exceções de erro distintas para caractere inválido e \emph{token}
  inesperado, que sustentam a separação entre erro léxico e sintático;
\item
  o \emph{parser} interativo, que permite alimentar \emph{tokens}
  manualmente e, assim, implementar a recuperação de erros em modo
  pânico.
\end{itemize}

A gramática reconhece apenas a forma das construções, isto é, a ordem das
seções, o formato de um gatilho, a estrutura de uma chamada, enquanto
regras que dependem de tipos, domínios e nomes ficam a cargo do
analisador semântico, que tem mais contexto para produzir mensagens
precisas.

\section{Definição da Linguagem}

O projeto da linguagem foi orientado por dois princípios fundamentais. O primeiro é fornecer uma sintaxe declarativa e uniforme, estruturada em seções com responsabilidades bem definidas. O segundo é adotar um vocabulário reduzido e composto por nomes amigáveis, como \texttt{turn\_on}, \texttt{delay} e \texttt{say}, cujo significado é conhecido e validado pelo compilador. Esta seção apresenta, em nível conceitual, as principais famílias de construções oferecidas pela linguagem.

\subsection{Automação}

A unidade básica de um programa é a automação, introduzida pela
palavra-chave \texttt{automation} seguida de um título e de um corpo
entre chaves. Um arquivo contém zero ou mais automações e,
opcionalmente, definições de nomes. O corpo é composto
por seções em ordem fixa: metadados opcionais
(\texttt{description}, \texttt{mode}, \texttt{max}), a seção
\texttt{when} (gatilhos, obrigatória), a seção \texttt{if} (condições,
opcional) e a seção \texttt{do} (ações, obrigatória). A
Listagem~\ref{lst:skeleton} apresenta o esqueleto de uma automação com
as quatro seções. A ordem metadados, \texttt{when}, \texttt{if} e
\texttt{do} é imposta pela gramática.

\begin{listing}[H]
    \centering
    \begin{minted}{swift}
automation "Aquecedor da suíte" {
    mode = "single"                              // metadado
    when (switch.toalheiro changes to "on")      // gatilho (obrigatório)
    if (between @15:00:00 @23:00:00)             // condição (opcional)
    do { turn_on(switch.toalheiro) }             // ações (obrigatório)
}
    \end{minted}
    \caption{Esqueleto de uma automação com as quatro seções.}
    \label{lst:skeleton}
\end{listing}

\subsection{Gatilhos}

A seção \texttt{when} lista um ou mais gatilhos, os eventos
que disparam a automação, separados pela palavra \texttt{or}. Cada
gatilho pode receber um identificador com \texttt{as\ "id"},
referenciado depois nas condições. Há quatro famílias:

\begin{itemize}
\item \textbf{de estado:} \texttt{ENTIDADE changes [from V] [to V] [for DURAÇÃO]}
  reage à mudança de estado de uma entidade, opcionalmente filtrando
  origem/destino e exigindo permanência por um tempo;
\item \textbf{numérico:} \texttt{ENTIDADE (\textless{} \textbar{} \textgreater{}) NÚMERO [for DURAÇÃO]}
  reage ao cruzamento de um limiar;
\item \textbf{de horário:} \texttt{@HH:MM:SS} dispara em um horário do
  dia;
\item \textbf{solar:} \texttt{(sunset \textbar{} sunrise) [(+ \textbar{} -) DURAÇÃO]}
  dispara no nascer ou pôr do sol, com deslocamento opcional.
\end{itemize}

A Listagem~\ref{lst:triggers} reúne as quatro famílias de gatilho, cada
uma com seu identificador.

\begin{listing}[H]
    \centering
    \begin{minted}{swift}
when (
    binary_sensor.sala changes to "on" for 5m as "movimento" or
    sensor.temperatura > 30 as "calor" or
    @07:00:00 as "manha" or
    sunset - 30m as "anoitecer"
)
    \end{minted}
    \caption{As quatro famílias de gatilho, cada uma com um identificador.}
    \label{lst:triggers}
\end{listing}

\subsection{Condições}

A seção \texttt{if} é uma expressão booleana que precisa ser
verdadeira para a automação prosseguir. As condições abrangem:

\begin{itemize}
  \item \textbf{estado:} \texttt{ENTIDADE == "v"} ou \texttt{!=};
  \item \textbf{numérica:} \texttt{ENTIDADE \textless{} N} ou \texttt{\textgreater{} N};
  \item \textbf{horário:} \texttt{after @T}, \texttt{before @T}, \texttt{between @T @T}, \texttt{on [dias]};
  \item \textbf{solar:} \texttt{after sunset}, \texttt{before sunrise -10m}, \texttt{between sunset sunrise} com deslocamento opcional;
  \item \textbf{gatilho:} \texttt{triggeredby "id"}.
\end{itemize}

Há suporte também para os operadores lógicos \texttt{not},
\texttt{and} e \texttt{or}, com precedência idêntica à linguagem
Python (\texttt{not} \textgreater{} \texttt{and} \textgreater{}
\texttt{or}) e parênteses para agrupar.

A Listagem~\ref{lst:conditions} combina uma comparação de estado, uma
comparação numérica, uma janela de horário e o predicado
\texttt{triggeredby} sob operadores lógicos.

\begin{listing}[H]
    \centering
    \begin{minted}{swift}
if (
    light.sala == "on" and
    sensor.temperatura < 28 and
    between @06:00:00 @23:00:00 and
    not triggeredby "calor"
)
    \end{minted}
    \caption{Condição combinando diferentes expressões.}
    \label{lst:conditions}
\end{listing}

A forma \texttt{between A B}, diferentemente de escrever
\texttt{after A and before B} como dois átomos, gera uma única
condição com ambos os limites, o que o Home Assistant interpreta como um
intervalo que pode atravessar dias distintos.

\subsection{Ações}

A seção \texttt{do} é uma sequência de comandos, executados em
ordem. As formas são chamadas de ação com \texttt{nome(arg, ...)}, esperas com \texttt{delay(DURAÇÃO)}, condicionais com \texttt{if}, \texttt{elif} e \texttt{else}, e blocos que continuam em caso de erro com \texttt{safe}.

A Listagem~\ref{lst:actions} exemplifica as formas de comando dentro de
\texttt{do}.

\begin{listing}[H]
    \centering
    \begin{minted}{swift}
do {
    turn_on(lampada, color=quente, transition=2s)   
    delay(5m)                                   
    if (triggeredby "manha") {
        turn_off(switch.luz)
    } else {
        toggle(switch.luz)
    }

    safe {
      say(tts.google, message="Aviso")
    }
}
    \end{minted}
    \caption{Formas de comando dentro de \texttt{do}.}
    \label{lst:actions}
\end{listing}

\subsection{Vocabulário de ações e invocações diretas}

As chamadas a serviços representam a parte mais importante de uma automação. No \texttt{jass}, uma ação depende diretamente do domínio do alvo, modificando o seu significado e os argumentos aceitos.
Isso significa que a mesma ação sobre dispositivos diferentes apresenta capacidades distintas:
\texttt{turn\_on} sobre uma luz aceita \texttt{color}; sobre um
interruptor, não. Pedir um argumento que a ação não oferece para aquele
domínio, usá-la em um domínio que ela não suporta, ou passar um valor de
tipo errado são erros de compilação. Os argumentos
posicionais são parâmetros nomeados implícitos, ligados pela ordem, e
precisam preceder os nomeados.

Por fim, como nenhum vocabulário fechado cobre todos os casos, o \texttt{jass}
permite invocar diretamente serviços com
\texttt{call("dominio.servico",\ ...)}, sendo a única ação isenta de validação de
argumentos.

A Listagem~\ref{lst:curated} demonstra as ações.

\begin{listing}[H]
    \centering
    \begin{minted}{swift}
// ok: luz aceita 'color'
turn_on(light.sala, color=rgb(255, 0, 0))   
// erro: interruptor não tem 'color'
// turn_on(switch.porch, color=rgb(255,0,0)) 
// erro: 'open' não suporta luzes
// open(light.sala)                          
call("vacuum.send_command", target=vacuum.roomba, data={"command": "start"})
    \end{minted}
    \caption{Ações dependentes de domínio e invocação direta de serviço.}
    \label{lst:curated}
\end{listing}

\subsection{Valores}

Os valores escritos em argumentos, metadados e nomes pertencem a um
conjunto fixo de tipos, resumido na Tabela~\ref{tab:value_types}. Tipos como
duração, horário, referência de entidade e
cor têm sintaxe própria.

\begin{table}[H]
\centering
\begin{tabular}{l l l}
\toprule
\textbf{Tipo} & \textbf{Exemplo} \\
\midrule
string & \texttt{"on"} \\
número & \texttt{30}, \texttt{0.09} \\
booleano & \texttt{True}, \texttt{False} \\
duração & \texttt{5m}, \texttt{1h30m}, \texttt{45s} \\
horário & \texttt{@07:00:00} \\
referência de entidade & \texttt{light.sala} \\
lista & \texttt{["off", "on"]} \\
dicionário & \texttt{\{"command": "start"\}} \\
cor & \texttt{rgb(255, 180, 80)} \\
\bottomrule
\end{tabular}
\caption{Tipos de valor da linguagem.}
\label{tab:value_types}
\end{table}

\subsection{Nomes e variáveis}

Para evitar repetição, o \texttt{jass}
permite vincular nomes a valores com
\texttt{def NOME = VALOR}. Esses nomes são associações de tempo
de compilação, existindo apenas durante a compilação e sendo substituídos
pelo valor correspondente na saída. Cada uso é verificado por tipo exatamente como se o valor
estivesse escrito no lugar.

Há dois escopos: o de arquivo (global, visível a
qualquer automação declarada abaixo) e o de bloco (local ao
\texttt{\{ \}} em que o nome é declarado, com sombreamento, isto é,
um nome local oculta o de mesmo nome em escopo externo.

A Listagem~\ref{lst:names} ilustra definições globais (\texttt{lampada},
\texttt{quente}) e local (\texttt{reforco}).

\begin{listing}[H]
    \centering
    \begin{minted}{swift}
def lampada = light.sala
def quente  = rgb(255, 180, 80)

automation "Luz aconchegante" {
    when (lampada changes to "on")
    do {
        turn_on(lampada, color=quente)
        def reforco = 100
        turn_on(lampada, brightness=reforco)
    }
}
    \end{minted}
    \caption{Definições globais e locais.}
    \label{lst:names}
\end{listing}

\section{Gramática Livre de Contexto (GLC)}

Esta seção formaliza a gramática livre
de contexto do \texttt{jass}, usando a notação EBNF (\emph{Extended
Backus--Naur Form}), com os seguintes símbolos:
\begin{itemize}
  \item \texttt{->} define uma produção;
  \item \texttt{\textbar{}} separa alternativas;
  \item \texttt{[\ x\ ]} indica um elemento opcional (zero ou uma ocorrência);
  \item \texttt{\{\ x\ \}} indica repetição (zero ou mais);
  \item \texttt{(\ ...\ )} agrupa.

\end{itemize}

Os terminais aparecem como literais
entre aspas (\texttt{\textquotesingle{}when\textquotesingle{}},
\texttt{\textquotesingle{}==\textquotesingle{}}) para palavras-chave e
pontuação, ou em maíusculas para as classes de \emph{token} reconhecidas
por expressão regular no analisador léxico. Os não-terminais aparecem em
minúsculas, como \texttt{automation}, \texttt{state\_trigger}.

% \subsection{Terminais e não-terminais}

A Tabela~\ref{tab:terminals} enumera os não-terminais e a Listagem~\ref{lst:grammar} reúne as produções da gramática.

% \begin{table}[H]
% \centering
% \begin{tabular}{l p{6.1cm} p{3.6cm}}
% \toprule
% \textbf{Classe de terminal} & \textbf{Exemplos / forma} & \textbf{Papel} \\
% \midrule
% Palavras-chave & \texttt{automation}, \texttt{when}, \texttt{if},
% \texttt{elif}, \texttt{else}, \texttt{do}, \texttt{safe}, \texttt{def},
% \texttt{changes}, \texttt{from}, \texttt{to}, \texttt{for}, \texttt{as},
% \texttt{or}, \texttt{and}, \texttt{not}, \texttt{after}, \texttt{before},
% \texttt{between}, \texttt{on}, \texttt{sunset}, \texttt{sunrise},
% \texttt{triggeredby}, \texttt{True}, \texttt{False}, \texttt{rgb} &
% lexemas fixos da linguagem \\
% \addlinespace
% Operadores e pontuação &
% \texttt{==\ !=\ =\ \textless\ \textgreater\ +\ -\ (\ )\ \{\ \}\ [\ ]\ ,\ :} &
% comparação, atribuição, deslocamento, agrupamento, separação \\
% \addlinespace
% \texttt{NAME} & identificador & nomes de ações, parâmetros e aliases \\
% \texttt{STRING} & \texttt{"texto"} & literal de texto \\
% \texttt{NUMBER} & \texttt{30}, \texttt{0.09} & literal numérico \\
% \texttt{DURATION} & \texttt{1h30m}, \texttt{45s} & literal de duração \\
% \texttt{TIME} & \texttt{@07:00:00} & literal de horário \\
% \texttt{ENTITY} & \texttt{light.sala} & referência de entidade \\
% \bottomrule
% \end{tabular}
% \caption{Classes de terminais da linguagem.}
% \label{tab:terminals}
% \end{table}

% \begin{table}[H]
% \centering
% \begin{tabular}{l p{6.1cm} p{3.6cm}}
% \toprule
% \textbf{Classe de terminal} & \textbf{Forma} & \textbf{Exemplos} \\
% \midrule
% Palavras-chave & \texttt{automation}, \texttt{when}, \texttt{if},
% \texttt{elif}, \texttt{else}, \texttt{do}, \texttt{safe}, \texttt{def},
% \texttt{changes}, \texttt{from}, \texttt{to}, \texttt{for}, \texttt{as},
% \texttt{or}, \texttt{and}, \texttt{not}, \texttt{after}, \texttt{before},
% \texttt{between}, \texttt{on}, \texttt{sunset}, \texttt{sunrise},
% \texttt{triggeredby}, \texttt{True}, \texttt{False}, \texttt{rgb} &
%  \\
% \addlinespace
% Operadores e pontuação &
% \texttt{==\ !=\ =\ \textless\ \textgreater\ +\ -\ (\ )\ \{\ \}\ [\ ]\ ,\ :} &
%  \\
% \addlinespace
% \texttt{NAME} & \texttt{/[a-zA-Z\_][a-zA-Z0-9\_]*/} & luz \\
% \texttt{STRING} & \texttt{/"(\textbackslash{}\textbackslash{}.\textbar{}[\textasciicircum{}"\textbackslash{}\textbackslash{}])*"/}   & \texttt{"texto"} \\
% \texttt{NUMBER} &  \texttt{/\textbackslash{}d+\textbackslash{}.\textbackslash{}d+\textbar{}\textbackslash{}d+/} & \texttt{30}, \texttt{0.09} \\
% \texttt{DURATION} & \texttt{/\textbackslash{}d+h(\textbackslash{}d+m)?(\textbackslash{}d+s)?\textbar{}\textbackslash{}d+m(\textbackslash{}d+s)?\textbar{}\textbackslash{}d+s/} & \texttt{1h30m}, \texttt{45s} \\
% \texttt{TIME} &  \texttt{/@\textbackslash{}d\{2\}:\textbackslash{}d\{2\}:\textbackslash{}d\{2\}/} &\texttt{@07:00:00} \\
% \texttt{ENTITY} & \texttt{/[a-z][a-z\_]*\textbackslash{}.[a-z0-9\_]+/} &\texttt{light.sala}  \\
% \bottomrule
% \end{tabular}
% \caption{Classes de terminais da linguagem.}
% \label{tab:terminals}
% \end{table}

\begin{table}[H]
\centering
\begin{tabular}{l p{6.1cm} p{3.6cm}}
\toprule
\textbf{Classe de terminal} & \textbf{Forma} \\
\midrule
Palavras-chave & \texttt{automation}, \texttt{when}, \texttt{if},
\texttt{elif}, \texttt{else}, \texttt{do}, \texttt{safe}, \texttt{def},
\texttt{changes}, \texttt{from}, \texttt{to}, \texttt{for}, \texttt{as},
\texttt{or}, \texttt{and}, \texttt{not}, \texttt{after}, \texttt{before},
\texttt{between}, \texttt{on}, \texttt{sunset}, \texttt{sunrise},
\texttt{triggeredby}, \texttt{True}, \texttt{False}, \texttt{rgb}
 \\
Operadores e pontuação &
\texttt{==\ !=\ =\ \textless\ \textgreater\ +\ -\ (\ )\ \{\ \}\ [\ ]\ ,\ :}
 \\
\addlinespace
\texttt{NAME} & \texttt{/[a-zA-Z\_][a-zA-Z0-9\_]*/} \\
\texttt{STRING} & \texttt{/"(\textbackslash{}\textbackslash{}.\textbar{}[\textasciicircum{}"\textbackslash{}\textbackslash{}])*"/}    \\
\texttt{NUMBER} &  \texttt{/\textbackslash{}d+\textbackslash{}.\textbackslash{}d+\textbar{}\textbackslash{}d+/}  \\
\texttt{DURATION} & \texttt{/\textbackslash{}d+h(\textbackslash{}d+m)?(\textbackslash{}d+s)?\textbar{}\textbackslash{}d+m(\textbackslash{}d+s)?\textbar{}\textbackslash{}d+s/} \\
\texttt{TIME} &  \texttt{/@\textbackslash{}d\{2\}:\textbackslash{}d\{2\}:\textbackslash{}d\{2\}/}  \\
\texttt{ENTITY} & \texttt{/[a-z][a-z\_]*\textbackslash{}.[a-z0-9\_]+/}  \\
\texttt{COMMENT} & \texttt{/\textbackslash{}/\textbackslash{}/[\textasciicircum{}\textbackslash{}n]*/} \\
\bottomrule
\end{tabular}
\caption{Classes de terminais da linguagem.}
\label{tab:terminals}
\end{table}

% TIME:     \texttt{/@\textbackslash{}d\{2\}:\textbackslash{}d\{2\}:\textbackslash{}d\{2\}/}
% DURATION: \texttt{/\textbackslash{}d+h(\textbackslash{}d+m)?(\textbackslash{}d+s)?\textbar{}\textbackslash{}d+m(\textbackslash{}d+s)?\textbar{}\textbackslash{}d+s/}
% NUMBER:   \texttt{/\textbackslash{}d+\textbackslash{}.\textbackslash{}d+\textbar{}\textbackslash{}d+/}
% ENTITY:   \texttt{/[a-z][a-z\_]*\textbackslash{}.[a-z0-9\_]+/}
% NAME:     \texttt{/[a-zA-Z\_][a-zA-Z0-9\_]*/}
% STRING:   \texttt{/"(\textbackslash{}\textbackslash{}.\textbar{}[\textasciicircum{}"\textbackslash{}\textbackslash{}])*"/}

% \begin{table}[H]
% \centering
% \begin{tabular}{p{6.4cm} p{6.4cm}}
% \toprule
% \textbf{Não-terminal(is)} & \textbf{Papel} \\
% \midrule
% \texttt{start} & programa: sequência de itens de topo \\
% \addlinespace
% \texttt{automation}, \texttt{metadata} & unidade de automação, suas
% seções e metadados \\
% \addlinespace
% \texttt{when\_block}, \texttt{trigger}, \texttt{state\_trigger},
% \texttt{numeric\_trigger}, \texttt{time\_trigger}, \texttt{sun\_trigger},
% \texttt{sun\_event} & gatilhos \\
% \addlinespace
% \texttt{condition\_block}, \texttt{bool\_expr}, \texttt{or\_expr},
% \texttt{and\_expr}, \texttt{not\_expr}, \texttt{atom} & expressão de
% condição (cascata de precedência) \\
% \addlinespace
% \texttt{condition} e os átomos \texttt{state\_condition},
% \texttt{numeric\_condition}, \texttt{time\_condition},
% \texttt{day\_condition}, \texttt{sun\_condition},
% \texttt{triggeredby\_condition} & átomos de condição \\
% \addlinespace
% \texttt{do\_block}, \texttt{block}, \texttt{statement} & corpo de ações \\
% \addlinespace
% \texttt{action\_call}, \texttt{arguments}, \texttt{pos\_args},
% \texttt{kw\_args}, \texttt{named\_arg} & chamada de ação e seus
% argumentos \\
% \addlinespace
% \texttt{if\_statement}, \texttt{elif\_clause}, \texttt{else\_clause},
% \texttt{safe\_block} & comandos compostos \\
% \addlinespace
% \texttt{def\_binding} & vínculo de nome (alias) \\
% \addlinespace
% \texttt{value}, \texttt{entity}, \texttt{threshold}, \texttt{list},
% \texttt{dict}, \texttt{pair}, \texttt{color} & valores \\
% \bottomrule
% \end{tabular}
% \caption{Principais não-terminais e seus papéis.}
% \label{tab:nonterminals}
% \end{table}

% \subsection{Produções e não-terminais}

% A Listagem~\ref{lst:grammar} reúne as produções da gramática do
% \texttt{jass}.

\begin{listing}[H]
    \centering
    \begin{minted}[fontsize=\scriptsize]{text}
// Programa e automação
start            -> { def_binding | automation }
def_binding      -> 'def' NAME '=' value
automation       -> 'automation' [ STRING ] '{' { metadata }
                       when_block [ condition_block ] do_block '}'
metadata         -> NAME '=' value

// Gatilhos
when_block       -> 'when' '(' trigger { 'or' trigger } ')'
trigger          -> trigger_expr [ 'as' STRING ]
trigger_expr     -> state_trigger | numeric_trigger | time_trigger | sun_trigger
state_trigger    -> entity 'changes' [ 'from' value ] [ 'to' value ] [ 'for' DURATION ]
numeric_trigger  -> entity ( '<' | '>' ) threshold [ 'for' DURATION ]
time_trigger     -> TIME
sun_trigger      -> sun_event
sun_event        -> ( 'sunset' | 'sunrise' ) [ ( '+' | '-' ) DURATION ]

// Condições (precedência: not > and > or)
condition_block  -> 'if' '(' bool_expr ')'
bool_expr        -> or_expr
or_expr          -> and_expr { 'or' and_expr }
and_expr         -> not_expr { 'and' not_expr }
not_expr         -> 'not' not_expr | atom
atom             -> '(' bool_expr ')' | condition
condition        -> state_condition | numeric_condition | time_condition
                 | day_condition | sun_condition | triggeredby_condition
state_condition  -> entity ( '==' | '!=' ) value
numeric_condition-> entity ( '<' | '>' ) threshold
time_condition   -> ( 'after' | 'before' ) TIME | 'between' TIME TIME
day_condition    -> 'on' list
sun_condition    -> ( 'after' | 'before' ) sun_event | 'between' sun_event sun_event
triggeredby_condition -> 'triggeredby' STRING
threshold        -> NUMBER | NAME

// Ações
do_block         -> 'do' block
block            -> '{' { statement } '}'
statement        -> action_call | if_statement | safe_block | def_binding
action_call      -> NAME '(' [ arguments ] ')'
arguments        -> pos_args | pos_args ',' kw_args | kw_args
pos_args         -> value | pos_args ',' value
kw_args          -> named_arg | kw_args ',' named_arg
named_arg        -> NAME '=' value
if_statement     -> 'if' '(' bool_expr ')' block { elif_clause } [ else_clause ]
elif_clause      -> 'elif' '(' bool_expr ')' block
else_clause      -> 'else' block
safe_block       -> 'safe' block

// Valores
entity           -> ENTITY | NAME
value            -> STRING | NUMBER | 'True' | 'False' | DURATION | TIME
                 | ENTITY | NAME | list | dict | color
list             -> '[' [ value { ',' value } ] ']'
dict             -> '{' [ pair { ',' pair } ] '}'
pair             -> STRING ':' value
color            -> 'rgb' '(' NUMBER ',' NUMBER ',' NUMBER ')'
    \end{minted}
    \caption{Produções da gramática do \texttt{jass}.}
    \label{lst:grammar}
\end{listing}

A gramática pode ser vista na íntegra no arquivo \href{https://github.com/jaimeadf-ufsm/elc408-compiladores/blob/main/trabalhos/t1/src/jass/grammar.lark}{grammar.lark}.

% \subsection{Decisões de projeto}

% \textbf{Precedência e associatividade dos operadores lógicos.} A
% precedência é codificada por uma \textbf{cascata de não-terminais}:
% \texttt{or\_expr} é definido sobre \texttt{and\_expr}, que é definido
% sobre \texttt{not\_expr}, que se reduz a \texttt{atom}. Como \texttt{or}
% só aparece no nível mais externo e \texttt{and} em um nível mais
% interno, obtém-se \texttt{not} \textgreater{} \texttt{and}
% \textgreater{} \texttt{or}, como em Python. As repetições
% \texttt{\{\ \textquotesingle{}or\textquotesingle{}\ ...\ \}} e
% \texttt{\{\ \textquotesingle{}and\textquotesingle{}\ ...\ \}} dão
% associatividade à esquerda, e a alternativa
% \texttt{atom\ ->\ \textquotesingle{}(\textquotesingle{}\ bool\_expr\ \textquotesingle{})\textquotesingle{}}
% permite reagrupar com parênteses.

% \textbf{Imposição da ordem das seções.} A produção de
% \texttt{automation} fixa a sequência
% \texttt{\{\ metadata\ \}\ when\_block\ {[}\ condition\_block\ {]}\ do\_block}.
% Assim, a ordem metadados → \texttt{when} → \texttt{if} → \texttt{do} é
% garantida \textbf{estruturalmente}: trocar a ordem (por exemplo,
% escrever \texttt{if} antes de \texttt{when}) ou omitir uma seção
% obrigatória é um erro sintático, não uma verificação posterior. O mesmo
% vale para a ausência de \texttt{or} pendente: \texttt{when\_block} exige
% um \texttt{trigger} após cada \texttt{or}, de modo que um \texttt{or} no
% fim da lista não deriva.

% \textbf{Argumentos posicionais antes dos nomeados.} A produção de
% \texttt{arguments}
% (\texttt{pos\_args\ \textbar{}\ pos\_args\ \textquotesingle{},\textquotesingle{}\ kw\_args\ \textbar{}\ kw\_args})
% só permite argumentos nomeados \textbf{depois} dos posicionais, nunca
% antes. Logo, \texttt{turn\_on(color=...,\ light.x)} é rejeitado já na
% análise sintática.

% \textbf{Classe de análise: LALR(1).} A gramática foi escrita para ser
% analisável por um \emph{parser} \textbf{LALR(1)} --- ascendente,
% dirigido por tabela e com \textbf{um} \emph{token} de \emph{lookahead}
% (símbolo de antecipação). Essa escolha favorece velocidade e produz
% pontos de erro precisos (um único \emph{token} inesperado). A gramática
% é projetada para que toda decisão se resolva com um símbolo de
% antecipação: após uma \texttt{entity}, o operador seguinte
% (\texttt{==}/\texttt{!=} versus
% \texttt{\textless{}}/\texttt{\textgreater{}}) distingue condição de
% estado de condição numérica; após \texttt{after}, o \emph{token}
% seguinte (\texttt{TIME} versus \texttt{sunset}/\texttt{sunrise})
% distingue janela de horário de janela solar. A construção do
% \emph{parser} não acusa conflitos. Por fim, os dois usos de \texttt{\{}
% \texttt{\}} (literal de dicionário e bloco de comandos) ocorrem em
% contextos disjuntos do \emph{parser}, de modo que o \textbf{lexer
% contextual} os reconhece sem ambiguidade.

\section{Análise Léxica}

A análise léxica converte o texto-fonte em uma sequência de
\emph{tokens}. No \texttt{jass}, essa fase é inteiramente delegada ao
Lark, sendo as definições de terminais declaradas na própria
gramática, das quais o Lark deriva o \emph{lexer}. Comentários de linha (\texttt{//\ ...}) e espaços em branco são
descartados por diretivas \texttt{\%ignore}, não chegando ao
\emph{parser}. 

% O Lark anota cada \emph{token} com sua linha e
% coluna de origem. Essa informação é propagada para a AST e é a base dos
% diagnósticos localizados.

% Esta seção descreve
%como esses terminais foram especificados e como os erros léxicos são
%detectados e relatados.

% \subsection{Especificação dos terminais no Lark}

% Cada terminal é declarado na gramática por uma \textbf{expressão
% regular} (para classes de \emph{token}) ou por um \textbf{literal de
% string} (para palavras-chave e pontuação). A
% Listagem~\ref{lst:terminals_regex} reproduz as definições das classes
% não triviais.

% \begin{listing}[H]
%     \centering
%     \begin{minted}{text}
% TIME:     /@\d{2}:\d{2}:\d{2}/
% DURATION: /\d+h(\d+m)?(\d+s)?|\d+m(\d+s)?|\d+s/
% NUMBER:   /\d+\.\d+|\d+/
% ENTITY:   /[a-z][a-z_]*\.[a-z0-9_]+/
% NAME:     /[a-zA-Z_][a-zA-Z0-9_]*/
% STRING:   /"(\\.|[^"\\])*"/
%     \end{minted}
%     \caption{Definições regulares das classes de \emph{token} mais
%     relevantes.}
%     \label{lst:terminals_regex}
% \end{listing}

% Três decisões merecem destaque, todas resolvidas pela combinação de
% \textbf{casamento mais longo} (\emph{longest match}) com a
% \textbf{prioridade} que o Lark dá a terminais literais sobre os
% definidos por expressão regular:

% \begin{itemize}
% \tightlist
% \item
%   \textbf{Palavras-chave versus identificadores.} As palavras-chave são
%   terminais literais (\texttt{"when"}, \texttt{"changes"}\ldots) e, por
%   isso, têm prioridade sobre \texttt{NAME}. Já um identificador que
%   apenas começa com uma palavra-chave (como \texttt{online}) é
%   reconhecido como \texttt{NAME}, pois o casamento mais longo prevalece.
% \item
%   \textbf{Durações versus números.} \texttt{DURATION} exige ao menos uma
%   unidade (\texttt{h}, \texttt{m} ou \texttt{s}); assim \texttt{45s}
%   casa como duração (lexema mais longo), enquanto \texttt{22} casa como
%   \texttt{NUMBER}. Em \texttt{NUMBER}, a alternativa de real
%   (\texttt{\textbackslash{}d+\textbackslash{}.\textbackslash{}d+})
%   precede a de inteiro pelo mesmo motivo.
% \item
%   \textbf{Referências de entidade.} \texttt{ENTITY} é um único
%   \emph{token} da forma \texttt{domínio.objeto}, em que o
%   \texttt{objeto} pode começar por dígito (p.ex.
%   \texttt{alarm\_control\_panel.37dc24}). Já \texttt{TIME} é capturado
%   como um \emph{token} único, o que mantém o \texttt{:} livre para uso
%   em dicionários.
% \end{itemize}



% \begin{listing}[H]
%     \centering
%     \begin{minted}{text}
% AUTOMATION  'automation'
% STRING      '"Sala"'
% LBRACE      '{'
% WHEN        'when'
% LPAR        '('
% ENTITY      'binary_sensor.sala'
% CHANGES     'changes'
% TO          'to'
% STRING      '"on"'
% FOR         'for'
% DURATION    '1h30m'
% RPAR        ')'
% ...
%     \end{minted}
%     \caption{\emph{Tokens} produzidos para
%     \texttt{when\ (binary\_sensor.sala\ changes\ to\ "on"\ for\ 1h30m)}.}
%     \label{lst:tokens_output}
% \end{listing}

% Internamente, o Lark compila cada definição regular em um autômato
% finito; a Figura~\ref{fig:duration_automaton} ilustra o reconhecimento de
% \texttt{DURATION}, a classe com a estrutura mais distintiva.

% \begin{figure}[htpb]
%     \centering
%     \begin{tikzpicture}[
%         node distance = 3.6cm,
%         on grid,
%         auto,
%         state/.style = {
%             circle, draw=black, thick,
%             minimum size=1cm, inner sep=1pt, align=center
%         },
%         accepting/.style = {double, double distance=1.2pt},
%         arrow/.style = {-{Stealth[scale=1.2]}, thick}
%     ]
%     \node[state] (q0) {$q_0$};
%     \node[state, accepting, right=of q0] (qh) {$q_h$};
%     \node[state, accepting, right=of qh] (qm) {$q_m$};
%     \node[state, accepting, right=of qm] (qf) {$q_f$};

%     \draw[arrow] ([xshift=-1cm]q0.west) -- (q0);
%     \draw[arrow] (q0) -- node {dígitos\,\texttt{h}} (qh);
%     \draw[arrow] (qh) -- node {dígitos\,\texttt{m}} (qm);
%     \draw[arrow] (qm) -- node {dígitos\,\texttt{s}} (qf);
%     \draw[arrow] (q0) to[bend left=38] node[above] {dígitos\,\texttt{m}} (qm);
%     \draw[arrow] (qh) to[bend left=38] node[above] {dígitos\,\texttt{s}} (qf);
%     \draw[arrow] (q0) to[bend right=32] node[below] {dígitos\,\texttt{s}} (qf);
%     \end{tikzpicture}
%     \caption{Esboço do autômato que reconhece \texttt{DURATION}
%     (\texttt{1h30m}, \texttt{45s}, \ldots): as unidades \texttt{h},
%     \texttt{m} e \texttt{s} aparecem nessa ordem, com ao menos uma
%     presente.}
%     \label{fig:duration_automaton}
% \end{figure}

% \subsection{Detecção e tratamento de erros léxicos}

Quando o \emph{lexer} encontra um caractere que
não inicia nenhum terminal, o Lark lança a exceção
\texttt{UnexpectedCharacters}, que o \texttt{jass} captura e
converte em um erro da categoria léxica, preservando a linha, a coluna e o caractere ofensor.
A Listagem~\ref{lst:lexical_error} mostra o
diagnóstico real para uma entrada que contém um \texttt{\$} inválido.

\begin{listing}[H]
    \centering
    \begin{minted}{text}
samples/errors/lexical.jass:5:27: lexical error: unexpected character '$'
          turn_on(light.sala$led)
                            ^
    \end{minted}
    \caption{Diagnóstico de erro léxico para o caractere \texttt{\$}.}
    \label{lst:lexical_error}
\end{listing}

% A separação clara entre essa categoria e as demais é intencional: o erro
% léxico é detectado antes da análise sintática e interrompe o
% processamento daquele arquivo com código de saída não nulo, sem ser
% confundido com um \emph{token} gramaticalmente inesperado (erro
% sintático, Seção 7).

\section{Análise Sintática}

A análise sintática verifica se a sequência de \emph{tokens} forma um
programa válido segundo a gramática e produz a árvore de
derivação. O \texttt{jass} usa o \emph{parser} LALR(1) do Lark, construído
a partir da gramática. As motivações já foram introduzidas (Seções 3 e 5):
velocidade, ausência de conflitos para esta gramática e pontos de erro
precisos. A análise opera em conjunto com o \textit{lexer} contextual,
que usa o estado do \emph{parser} para escolher, quando há ambiguidade
léxica, apenas entre os terminais válidos naquele ponto.

% As \textbf{regras estruturais} da linguagem são impostas pela própria
% gramática e, por isso, manifestam-se como erros sintáticos: a ordem fixa
% das seções, a proibição de \texttt{or} pendente e a exigência de
% argumentos posicionais antes dos nomeados (Subseção 5.3). A
% Listagem~\ref{lst:parse_tree} mostra a árvore de derivação real
% (\texttt{-\/-emit\ parse}) de uma automação mínima, na qual se
% reconhecem as seções \texttt{when\_block} e \texttt{do\_block} e o
% aninhamento até os átomos.

% \begin{listing}[H]
%     \centering
%     \begin{minted}{text}
% start
%   automation
%     "Sala"
%     when_block
%       trigger
%         state_trigger
%           entity_ref  light.sala
%           to
%           string  "on"
%     do_block
%       block
%         action_call
%           turn_off
%           arguments
%             pos_args
%               entity_ref  light.sala
%     \end{minted}
%     \caption{Árvore de derivação (simplificada, sem os \emph{tokens} de
%     pontuação) de \texttt{automation\ "Sala"\ \{\ when\ (light.sala\
%     changes\ to\ "on")\ do\ \{\ turn\_off(light.sala)\ \}\ \}}.}
%     \label{lst:parse_tree}
% \end{listing}

% A Figura~\ref{fig:condition_tree} detalha a derivação de uma
% \textbf{condição} que exercita a cascata de precedência:
% \texttt{light.a\ ==\ "on"\ and\ not\ light.b\ ==\ "off"}. A árvore
% evidencia que \texttt{and} está acima de \texttt{not} (logo,
% \texttt{not} liga-se mais fortemente) e que cada comparação é um
% \texttt{state\_condition}.

% \begin{figure}[htpb]
%     \centering
%     \resizebox{\textwidth}{!}{%
%     \begin{tikzpicture}[
%         treenode/.style = {
%             rectangle, draw=black, thick, rounded corners=3pt,
%             align=center, inner sep=4pt, font=\footnotesize
%         },
%         level distance = 1.6cm,
%         level 1/.style = {sibling distance=4.6cm},
%         level 2/.style = {sibling distance=4cm},
%         edge from parent/.style = {draw=black, thick, -{Stealth[scale=1.0]}},
%         every node/.style = treenode
%     ]
%     \node {bool\_expr $\to$ or\_expr $\to$ and\_expr}
%       child { node {not\_expr $\to$ atom}
%         child { node {state\_condition\\ \texttt{light.a == "on"}} }
%       }
%       child { node {\texttt{'and'}} }
%       child { node {not\_expr}
%         child { node {\texttt{'not'}} }
%         child { node {atom}
%           child { node {state\_condition\\ \texttt{light.b == "off"}} }
%         }
%       };
%     \end{tikzpicture}%
%     }
%     \caption{Árvore de derivação de
%     \texttt{light.a\ ==\ "on"\ and\ not\ light.b\ ==\ "off"}, mostrando
%     a precedência \texttt{not} \textgreater{} \texttt{and}.}
%     \label{fig:condition_tree}
% \end{figure}

% \subsection{Detecção e recuperação de erros}

Quando o \emph{parser} recebe um \emph{token} que nenhuma entrada da
tabela aceita, o Lark sinaliza um \emph{token} inesperado, que o
\texttt{jass} converte em um erro da categoria sintática, com posição e
a lista de \emph{tokens} esperados naquele ponto. Em vez de
parar no primeiro erro, o \texttt{jass} aplica
recuperação em modo pânico, isto é, ao detectar a falha, registra o
erro, descarta \emph{tokens} até alcançar um ponto de
sincronização e retoma a análise a partir dele. Os pontos de
sincronização são as \textbf{fronteiras de topo} do programa, as
palavras-chave \texttt{automation} e \texttt{def}, que marcam o
início de uma nova unidade independente. Assim, um erro em uma automação
não contamina as seguintes, e uma única execução coleta todos
os erros. A  Listagem~\ref{lst:syntactic_errors} mostra os erros gerados 
para um arquivo com três automações, cada uma com um
defeito distinto: bloco \texttt{do} ausente, \texttt{or} pendente
(esperava-se um novo gatilho após \texttt{or}) e seção \texttt{if} antes
de \texttt{when}.

% Para ter esse controle, o \texttt{jass} não usa o laço de análise
% padrão: ele alimenta os \emph{tokens} um a um no \textbf{\emph{parser}
% interativo} do Lark; ao ocorrer uma exceção, executa a ressincronização
% descrita e prossegue. A Listagem~\ref{lst:syntactic_errors} mostra o
% resultado real para um arquivo com três automações, cada uma com um
% defeito distinto: bloco \texttt{do} ausente, \texttt{or} pendente
% (esperava-se um novo gatilho após \texttt{or}) e seção \texttt{if} antes
% de \texttt{when}.

\begin{listing}[H]
    \centering
    \begin{minted}{text}
samples/errors/syntactic.jass:5:1:  syntactic error: unexpected '}'; expected 'do'
samples/errors/syntactic.jass:12:5: syntactic error: unexpected ')'; expected 'sunrise', 'sunset', a name, a time literal, ...
samples/errors/syntactic.jass:20:5: syntactic error: unexpected 'if'; expected 'when', a name
    \end{minted}
    \caption{Três erros sintáticos relatados em uma só execução. Após
    cada falha, o \emph{parser} avança até a próxima \texttt{automation}
    e recomeça.}
    \label{lst:syntactic_errors}
\end{listing}

% Concretamente, no primeiro defeito o \emph{parser} esperava a seção
% \texttt{do} e encontrou \texttt{\}}; ele descarta os \emph{tokens}
% restantes daquela automação até a palavra \texttt{automation} da segunda
% e retoma --- repetindo o processo até o fim do arquivo. Os erros léxicos
% e semânticos permanecem em categorias separadas: o léxico é detectado
% antes (Seção 6) e o semântico, depois (Seção 8).

\section{Análise Semântica}

A análise semântica é a fase em que o \texttt{jass} concentra a maior
parte de suas verificações, percorrendo a AST de cada automação e coletando 
todos os erros. A análise é baseada em três pilares, tratados a seguir: a tabela de
símbolos, a verificação de tipos e a consistência ação e domínio.

\subsection{Tabela de símbolos e escopos}

A tabela de símbolos registra os nomes vinculados por
\texttt{def}, associando a cada um o tipo inferido do valor e
o próprio valor. Os
tipos formam um conjunto fixo de categorias, como string, número,
booleano, duração, horário, entidade (que carrega também o domínio), lista, dicionário.

Os escopos são organizados em uma pilha aninhada que reflete a
estrutura do programa: o escopo de arquivo (global), o de cada
automação e o de cada bloco. A resolução de
um nome percorre a pilha do escopo atual para fora. Isso significa que uma definição local
de mesmo nome ``sombreia'' a externa. Nomes não resolvidos são erros. As referências dos
gatilhos declaradas com \texttt{as\ "id"} são coletadas por automação e
formam o conjunto contra o qual cada \texttt{triggeredby\ "id"} é
validado, inclusive nos \texttt{if} internos às ações.

%A
%Figura~\ref{fig:scopes} ilustra os escopos e o sombreamento.

% \begin{figure}[htpb]
%     \centering
%     \begin{tikzpicture}[
%         node distance = 1.0cm,
%         box/.style = {
%             rectangle, draw=black, thick, align=center,
%             rounded corners=3pt, minimum width=8.5cm,
%             minimum height=1cm, inner sep=5pt
%         },
%         arrow/.style = {-{Stealth[scale=1.2]}, thick}
%     ]
%     \node[box] (g) {\textbf{Escopo de arquivo (global)}\\
%         \texttt{lampada}: entidade(\texttt{light}) $\cdot$ \texttt{quente}: cor};
%     \node[box, below=of g] (a) {\textbf{Escopo da automação}\\
%         ids de gatilho: \{ \texttt{manha} \}};
%     \node[box, below=of a] (b) {\textbf{Escopo do bloco \texttt{do}}\\
%         \texttt{reforco}: número};
%     \node[box, below=of b] (i) {\textbf{Escopo do ramo \texttt{if}}\\
%         \texttt{quente}: cor (sombreia o global)};
%     \draw[arrow] (g) -- (a);
%     \draw[arrow] (a) -- (b);
%     \draw[arrow] (b) -- (i);
%     \end{tikzpicture}
%     \caption{Pilha de escopos. O \texttt{quente} declarado no ramo
%     \texttt{if} sombreia o global homônimo; \texttt{reforco} só existe
%     dentro do bloco \texttt{do}. A resolução vai do escopo interno ao
%     externo.}
%     \label{fig:scopes}
% \end{figure}

\subsection{Verificação de tipos}

Cada valor usado em um argumento, gatilho ou condição tem seu tipo
inferido e comparado ao tipo esperado naquele ponto. Quando o
valor é uma referência a uma definição, ela é resolvida ao seu
tipo antes de qualquer verificação. A Listagem~\ref{lst:type_check}
exemplifica esse tipo de erro. 



%A
% Listagem~\ref{lst:type_check} contrasta um uso válido com um quase
% idêntico que falha.

% um alias de duração onde se espera uma cor falha
% exatamente como se a duração estivesse escrita no lugar. Combinações
% inválidas --- um limiar numérico que não é número, um argumento de cor
% que recebe uma duração --- são reportadas com localização. A
% Listagem~\ref{lst:type_check} contrasta um uso válido com um quase
% idêntico que falha.

\begin{listing}[H]
    \centering
    \begin{minted}{swift}
// válido: 'color' recebe uma cor
turn_on(light.sala, color=rgb(255, 0, 0))   

def grace = 10m
// inválido
turn_on(light.sala, color=grace)            
    \end{minted}
    \begin{minted}{text}
semantic error: argument 'color' expects a color, got a duration
    do { turn_on(light.sala, color=grace) }
                                   ^^^^^
    \end{minted}
    \caption{Verificação de tipo.}
    \label{lst:type_check}
\end{listing}

\subsection{Consistência ação e domínio}

Este é o pilar central do vocabulário curado. O significado de uma ação
depende do \textbf{domínio do alvo} --- o tipo do dispositivo indicado
pelo primeiro parâmetro. A análise procede em etapas: (i) determina o
domínio do alvo, exigindo que, num alvo em lista, \textbf{todas as
entidades compartilhem o mesmo domínio}; (ii) verifica que a ação
\textbf{é definida para aquele domínio}; (iii) valida os argumentos ---
cada argumento precisa ser aceito por aquela ação \textbf{naquele
domínio}, com o tipo correto, respeitando obrigatoriedade, ausência de
duplicatas e a ligação posicional pela ordem. A única exceção é
\texttt{call(...)}, isenta de validação.

O conhecimento que sustenta essas verificações é o mesmo que guia a
tradução (Seção 10): para cada par (ação, domínio), o compilador conhece
os argumentos aceitos, seus tipos e como cada um é mapeado para o
serviço do Home Assistant. A Tabela~\ref{tab:turn_on_domains} mostra
como \texttt{turn\_on} difere entre dois domínios.

\begin{table}[H]
\centering
\begin{tabular}{l l p{6.8cm}}
\toprule
\textbf{Ação} & \textbf{Domínio} & \textbf{Argumentos aceitos} \\
\midrule
\texttt{turn\_on} & \texttt{light} & \texttt{target} (entidade),
\texttt{color} (cor), \texttt{brightness} (número), \texttt{transition}
(duração) \\
\addlinespace
\texttt{turn\_on} & \texttt{switch} & \texttt{target} (entidade) \\
\bottomrule
\end{tabular}
\caption{A mesma ação expõe argumentos diferentes conforme o domínio do
alvo: \texttt{color} é válido em \texttt{light} e inexistente em
\texttt{switch}.}
\label{tab:turn_on_domains}
\end{table}


A Listagem~\ref{lst:domain_errors} exibe as três classes de falha desta
fase --- argumento ausente para o domínio, ação incompatível com o
domínio e lista de alvos heterogênea ---, com os diagnósticos reais.

\begin{listing}[H]
    \centering
    \begin{minted}{text}
// (a) argumento inexistente para o domínio
turn_on(switch.porch, color=rgb(255, 0, 0))
  => action 'turn_on' has no argument 'color'

// (b) ação não suportada pelo domínio
open(light.sala)
  => action 'open' does not support domain 'light' (supported: cover)

// (c) domínios misturados no alvo
turn_on([light.sala, switch.ventilador])
  => all entities in a target list must share one domain (found: light, switch)
    \end{minted}
    \caption{Falhas de consistência ação $\leftrightarrow$ domínio.}
    \label{lst:domain_errors}
\end{listing}


\section{Detecção e Tratamento de
Erros}

Um objetivo central do \texttt{jass} é que as falhas apareçam \textbf{na
própria linguagem}, antes de qualquer YAML ser gerado. Esta seção
consolida a história de erros já distribuída pelas Seções 5 a 7,
padronizando como cada categoria é detectada, relatada e recuperada.

\subsection{As três categorias}

Os erros pertencem a três categorias, cada uma detectada em uma fase
distinta do \emph{pipeline}. A Tabela 6 resume o panorama; os mecanismos
de detecção foram detalhados nas seções por fase e não são repetidos
aqui.


As fases são ordenadas, e um erro em uma delas \textbf{impede} o avanço
para a seguinte: um arquivo com erro léxico nunca chega ao
\emph{parser}, e a análise semântica só ocorre se a sintática não acusou
falhas. Em consequência, uma execução relata erros de \textbf{uma única}
categoria --- a da primeira fase que falhou --- evitando diagnósticos
derivados e confusos.

\subsection{Formato e coleta dos
diagnósticos}

Internamente, as três categorias são representadas por uma hierarquia de
exceções (\texttt{LexicalError}, \texttt{SyntacticError},
\texttt{SemanticError}), cada uma carregando uma \textbf{localização}
(linha e coluna). Todas são impressas em um \textbf{formato uniforme}
---
\texttt{arquivo:linha:coluna:\ \textless{}categoria\textgreater{}\ error:\ \textless{}mensagem\textgreater{}}
--- seguido da linha de origem e de um marcador \texttt{\^{}} sob o
ponto da falha. Em caso de erro, o compilador termina com \textbf{código
de saída não nulo} (65); um arquivo válido sai com 0.

Quanto à coleta, há uma diferença deliberada: o erro \textbf{léxico}
interrompe o arquivo no primeiro caractere inválido, ao passo que os
erros \textbf{sintáticos} e \textbf{semânticos} são \textbf{acumulados}
--- pela recuperação em modo pânico (§6.2) e pela varredura completa da
AST (§7), respectivamente ---, de modo que vários problemas da mesma
categoria apareçam em uma só execução. O Trecho 15 reúne um exemplo
representativo por categoria, com o diagnóstico real emitido.


\section{Geração de Código}

A geração de código traduz a AST validada para a estrutura de saída do
Home Assistant. O processo tem dois passos: primeiro a AST é convertida
em um \textbf{modelo de saída} --- uma árvore de dicionários e listas
Python já no formato da plataforma ---, que em seguida é
\textbf{serializado} em YAML. Esta seção descreve o mapeamento de cada
família de construção, a resolução das ações curadas e o tratamento da
formatação.

\subsection{Estrutura geral}

Cada automação vira um dicionário com as chaves, em ordem fixa,
\texttt{alias}, \texttt{description} (opcional), \texttt{mode} (padrão
\texttt{single}), \texttt{max} (opcional), \texttt{triggers},
\texttt{conditions} (opcional) e \texttt{actions}; o arquivo inteiro
torna-se uma \textbf{lista} YAML de automações. A
Tabela~\ref{tab:codegen_mapping} resume o mapeamento por família de
construção.

\begin{table}[H]
\centering
\begin{tabular}{p{5.5cm} p{7.3cm}}
\toprule
\textbf{Construção \texttt{jass}} & \textbf{Construção de saída (Home Assistant)} \\
\midrule
gatilho de estado / numérico / horário / solar &
\texttt{trigger:\ state} / \texttt{numeric\_state} / \texttt{time} /
\texttt{sun} \\
\addlinespace
\texttt{changes\ from/to/for}, limiar \texttt{\textless}/\texttt{\textgreater} &
\texttt{from}/\texttt{to}/\texttt{for}, \texttt{below}/\texttt{above} \\
\addlinespace
\texttt{and} / \texttt{or} / \texttt{not} & formas abreviadas
\texttt{and:} / \texttt{or:} / \texttt{not:} \\
\addlinespace
\texttt{==} / \texttt{!=} & \texttt{condition:\ state} /
\texttt{condition:\ state} dentro de \texttt{not:} \\
\addlinespace
\texttt{after}/\texttt{before}/\texttt{between} (horário, sol) &
\texttt{condition:\ time}/\texttt{sun} com \texttt{after}/\texttt{before} \\
\addlinespace
\texttt{on\ [dias]} & \texttt{condition:\ time} com \texttt{weekday} \\
\addlinespace
\texttt{triggeredby\ "id"} & \texttt{condition:\ trigger} com \texttt{id} \\
\addlinespace
chamada de ação curada & bloco \texttt{action:} com
\texttt{target}/\texttt{data} \\
\addlinespace
\texttt{delay(d)} & \texttt{delay:\ HH:MM:SS} \\
\addlinespace
\texttt{if/elif/else} & \texttt{if/then/else} (cadeia de \texttt{elif}
aninhada em \texttt{else}) \\
\addlinespace
\texttt{safe\ \{\ \}} & \texttt{sequence:} com
\texttt{continue\_on\_error:\ true} nos filhos \\
\bottomrule
\end{tabular}
\caption{Mapeamento das construções \texttt{jass} para o YAML do Home
Assistant.}
\label{tab:codegen_mapping}
\end{table}

\subsection{Resolução das ações curadas}

A tradução de uma ação \textbf{não é uma regra ingênua} do tipo
\texttt{domínio.ação}: o compilador consulta, para o par (ação,
domínio), a \textbf{mesma definição} usada na análise semântica (Seção
8), que indica o serviço concreto do Home Assistant e \textbf{como cada
argumento é roteado} --- para \texttt{target.entity\_id} ou para uma
chave sob \texttt{data} --- e \textbf{convertido} conforme seu tipo.
Assim, \texttt{turn\_on} sobre uma luz resolve-se para o serviço
\texttt{light.turn\_on}, \texttt{toggle} sobre um interruptor para
\texttt{switch.toggle}, \texttt{open} sobre uma cortina para
\texttt{cover.open\_cover}, e assim por diante. As conversões por tipo
incluem cores em listas \texttt{{[}r,\ g,\ b{]}}, durações em
\texttt{HH:MM:SS} (ou em segundos quando o campo assim o exige, como
\texttt{transition}) e a resolução de aliases ao valor vinculado. A
Listagem~\ref{lst:codegen_curated} mostra uma chamada curada e o bloco
que ela gera.

\begin{listing}[H]
    \centering
    \begin{minted}{swift}
turn_on(light.sala, color=rgb(255, 0, 0), transition=5s)
    \end{minted}
    \begin{minted}{yaml}
- action: light.turn_on
  target:
    entity_id: light.sala
  data:
    rgb_color:        # cor -> lista [r, g, b]
      - 255
      - 0
      - 0
    transition: 5     # duração -> segundos (campo 'transition')
    \end{minted}
    \caption{Resolução de uma ação curada: \texttt{turn\_on} no domínio
    \texttt{light} vira o serviço \texttt{light.turn\_on}, com cada
    argumento roteado e convertido.}
    \label{lst:codegen_curated}
\end{listing}

A \texttt{call(...)} é traduzida diretamente: o primeiro argumento vira
o serviço, e \texttt{target}/\texttt{data} são repassados como estão ---
coerente com seu papel de escotilha de escape.

\subsection{Comandos compostos e nomes}

Os comandos compostos preservam a estrutura aninhada. Um
\texttt{if/elif/else} vira o bloco \texttt{if/then/else} do Home
Assistant, com cada \texttt{elif} aninhado no \texttt{else} do anterior;
um \texttt{safe\ \{\ \}} vira um bloco \texttt{sequence} cujos filhos
recebem \texttt{continue\_on\_error:\ true}. Quanto aos \textbf{nomes},
os aliases são \textbf{resolvidos ao seu valor vinculado} durante a
geração, respeitando o sombreamento de escopos; como são construções de
tempo de compilação, \textbf{nada sobre eles aparece no YAML}. A
Listagem~\ref{lst:codegen_composite} reúne ambos os casos a partir de
uma automação real.

\begin{listing}[H]
    \centering
    \begin{minted}{yaml}
# if/elif/else  ->  if/then/else
- if:
    - condition: trigger
      id: "anoitecer"
  then:
    - action: switch.toggle
      target:
        entity_id: switch.luz
  else:
    - delay: "00:05:00"

# safe { }  ->  sequence com continue_on_error
- sequence:
    - action: tts.speak
      target:
        entity_id: tts.google
      data:
        message: "Aviso"
      continue_on_error: true
    \end{minted}
    \caption{Saída real para um condicional e um bloco \texttt{safe}.}
    \label{lst:codegen_composite}
\end{listing}

\subsection{Formatação}

A serialização usa a biblioteca \texttt{ruamel.yaml} em \textbf{estilo
de bloco}, com escalares de texto entre aspas duplas e \textbf{ordem de
chaves estável} (a ordem de inserção no modelo). Horários e durações são
emitidos como strings (\texttt{"-00:30:00"}, \texttt{"00:05:00"}),
evitando interpretações indevidas pelo carregador YAML. O resultado é um
documento pronto para ser colado na configuração de automações da
plataforma.

\section{Exemplos Completos}

Esta seção reúne automações \textbf{completas} em \texttt{jass} ao lado
de sua saída \textbf{completa}, demonstrando o pipeline de ponta a
ponta. São dois casos: o primeiro, percorrido em detalhe, exercita
gatilhos variados, condições com janela solar, o recurso de nomes e
ações curadas; o segundo, um exemplo real do conjunto de testes,
exercita gatilhos numéricos, uma cadeia \texttt{elif} e a escotilha de
escape.

\subsection{Exemplo detalhado: luz da sala ao anoitecer}

A Listagem~\ref{lst:example1_src} mostra a fonte. Ela define três
\textbf{aliases} globais e uma automação que acende a luz com cor
aconchegante quando há movimento ou ao anoitecer, desde que seja noite e
a casa esteja desarmada.

\begin{listing}[H]
    \centering
    \begin{minted}{swift}
def lampada   = light.sala
def aconchego = rgb(255, 150, 60)
def espera    = 5m

automation "Sala ao anoitecer" {
    description = "Luz aconchegante quando ha movimento ao anoitecer"
    mode = "restart"

    when (
        binary_sensor.sala changes to "on" as "movimento" or
        sunset - 15m as "anoitecer"
    )

    if (
        between sunset sunrise and
        alarm_control_panel.casa == "disarmed"
    )

    do {
        turn_on(lampada, color=aconchego, transition=2s)
        if (triggeredby "movimento") {
            delay(espera)
            turn_off(lampada)
        } else {
            say(tts.sala, message="Bom anoitecer")
        }
    }
}
    \end{minted}
    \caption{Fonte do exemplo detalhado.}
    \label{lst:example1_src}
\end{listing}

A Listagem~\ref{lst:example1_yaml} mostra o YAML gerado, na íntegra.

\begin{listing}[H]
    \centering
    \begin{minted}[fontsize=\scriptsize]{yaml}
- alias: "Sala ao anoitecer"
  description: "Luz aconchegante quando ha movimento ao anoitecer"
  mode: "restart"
  triggers:
    - trigger: state
      entity_id: binary_sensor.sala
      to: "on"
      id: "movimento"
    - trigger: sun
      event: sunset
      offset: "-00:15:00"
      id: "anoitecer"
  conditions:
    - condition: sun
      after: sunset
      before: sunrise
    - condition: state
      entity_id: alarm_control_panel.casa
      state: "disarmed"
  actions:
    - action: light.turn_on
      target:
        entity_id: light.sala
      data:
        rgb_color:
          - 255
          - 150
          - 60
        transition: 2
    - if:
        - condition: trigger
          id: "movimento"
      then:
        - delay: "00:05:00"
        - action: light.turn_off
          target:
            entity_id: light.sala
      else:
        - action: tts.speak
          target:
            entity_id: tts.sala
          data:
            message: "Bom anoitecer"
    \end{minted}
    \caption{Saída completa para a Listagem~\ref{lst:example1_src}.}
    \label{lst:example1_yaml}
\end{listing}

Percorrendo a tradução:

\begin{itemize}
\tightlist
\item
  \textbf{Metadados.} \texttt{description} e \texttt{mode} viram chaves
  homônimas; como há \texttt{mode} explícito (\texttt{restart}), ele
  substitui o padrão \texttt{single}.
\item
  \textbf{Gatilhos.} O gatilho de estado vira \texttt{trigger:\ state}
  com \texttt{to} e o identificador \texttt{movimento}; o gatilho solar
  vira \texttt{trigger:\ sun} com \texttt{event:\ sunset} e
  \texttt{offset:\ "-00:15:00"} --- note a conversão de \texttt{-\ 15m}
  para o deslocamento assinado em \texttt{HH:MM:SS}.
\item
  \textbf{Condições.} A janela \texttt{between\ sunset\ sunrise} gera
  uma \textbf{única} condição \texttt{sun} com \texttt{after} e
  \texttt{before}, permitindo o intervalo que cruza a meia-noite; a
  comparação de estado vira uma condição \texttt{state}.
\item
  \textbf{Ações e nomes.} Na primeira ação, os aliases são resolvidos:
  \texttt{lampada} → \texttt{light.sala}, \texttt{aconchego} → a lista
  \texttt{{[}255,\ 150,\ 60{]}} em \texttt{rgb\_color}, e
  \texttt{espera}, usado em \texttt{delay}, → \texttt{"00:05:00"}. A
  ação \texttt{turn\_on} no domínio \texttt{light} resolve-se para
  \texttt{light.turn\_on}, com \texttt{transition} convertido para
  segundos. O \texttt{if/else} interno vira \texttt{if/then/else}, e o
  \texttt{triggeredby\ "movimento"} vira uma condição \texttt{trigger}.
  Nenhum dos aliases aparece no YAML, confirmando seu caráter de tempo
  de compilação.
\end{itemize}

\subsection{Exemplo real: carregador do tablet}

A Listagem~\ref{lst:example2_src} traz uma automação real do conjunto de
testes, que gerencia o carregamento de um tablet por faixas de bateria,
e a Listagem~\ref{lst:example2_yaml}, o YAML correspondente. Ela exercita
\textbf{quatro gatilhos numéricos} com identificadores, uma cadeia
\texttt{if/elif} com \texttt{or} entre predicados \texttt{triggeredby}, e
a escotilha de escape \texttt{call} para o serviço de notificação.

\begin{listing}[H]
    \centering
    \begin{minted}{swift}
automation "Tablet - Carregador" {
    mode = "single"
    when (
        sensor.d46e6d < 20 as "- de 20" or
        sensor.d46e6d > 75 as "+ de 75" or
        sensor.d46e6d < 10 as "- de 10" or
        sensor.d46e6d < 5  as "- de 5"
    )
    do {
        if (triggeredby "+ de 75") {
            turn_off(switch.6a55b6)
            call("notify.mobile_app_zfold4", data={"message": "Tablet descarregando"})
        } elif (triggeredby "- de 20" or triggeredby "- de 10") {
            turn_on(switch.6a55b6)
            call("notify.mobile_app_zfold4", data={"message": "Tablet carregando"})
        } elif (triggeredby "- de 5") {
            turn_on(switch.6a55b6)
            call("notify.mobile_app_zfold4",
                 data={"message": "Tablet carregando", "title": "Tablet zerado de bateria"})
        }
    }
}
    \end{minted}
    \caption{Fonte do exemplo real
    (\texttt{samples/automations/02-tablet-carregador.jass}).}
    \label{lst:example2_src}
\end{listing}

\begin{listing}[H]
    \centering
    \begin{minted}[fontsize=\scriptsize]{yaml}
- alias: "Tablet - Carregador"
  mode: "single"
  triggers:
    - trigger: numeric_state
      entity_id: sensor.d46e6d
      below: 20
      id: "- de 20"
    - trigger: numeric_state
      entity_id: sensor.d46e6d
      above: 75
      id: "+ de 75"
    - trigger: numeric_state
      entity_id: sensor.d46e6d
      below: 10
      id: "- de 10"
    - trigger: numeric_state
      entity_id: sensor.d46e6d
      below: 5
      id: "- de 5"
  actions:
    - if:
        - condition: trigger
          id: "+ de 75"
      then:
        - action: switch.turn_off
          target:
            entity_id: switch.6a55b6
        - action: "notify.mobile_app_zfold4"
          data:
            message: "Tablet descarregando"
      else:
        - if:
            - or:
                - condition: trigger
                  id: "- de 20"
                - condition: trigger
                  id: "- de 10"
          then:
            - action: switch.turn_on
              target:
                entity_id: switch.6a55b6
            - action: "notify.mobile_app_zfold4"
              data:
                message: "Tablet carregando"
          else:
            - if:
                - condition: trigger
                  id: "- de 5"
              then:
                - action: switch.turn_on
                  target:
                    entity_id: switch.6a55b6
                - action: "notify.mobile_app_zfold4"
                  data:
                    message: "Tablet carregando"
                    title: "Tablet zerado de bateria"
    \end{minted}
    \caption{Saída completa para a Listagem~\ref{lst:example2_src}. A
    cadeia \texttt{elif} aninha-se sucessivamente nos blocos
    \texttt{else}, e cada \texttt{call} é repassado diretamente como um
    bloco \texttt{action:} com seu \texttt{data}.}
    \label{lst:example2_yaml}
\end{listing}

Os dois exemplos, validados pelo compilador e carregáveis pela
plataforma, evidenciam que o pipeline percorre corretamente todas as
fases --- da análise léxica à geração de código --- para entradas que
exercitam a amplitude da linguagem.

\section{Conclusão}

O \texttt{jass} cumpre os objetivos propostos: oferece uma sintaxe
declarativa e legível que esconde o \emph{boilerplate} do YAML de
automações do Home Assistant e desloca a detecção de erros para o tempo
de compilação, relatando falhas léxicas, sintáticas e semânticas com
categoria e localização --- e, nas duas últimas, coletando vários
problemas em uma única execução --- antes de gerar qualquer saída; todas
as automações de exemplo são transpiladas para YAML válido e carregável,
validando o pipeline de ponta a ponta.

Como limitações conhecidas, o
vocabulário de ações é curado e cobre apenas os pares (ação, domínio)
implementados, de modo que ampliá-lo exige estender o registro de ações;
a política de domínios é leniente (domínios desconhecidos só são
rejeitados onde um serviço precisa ser mapeado) e a ação \texttt{call}
permanece, por projeto, sem validação de argumentos. Como extensões
plausíveis, destacam-se a ampliação do registro de ações e domínios, o
suporte a novas formas de gatilho e condição, a integração com a
validação nativa do Home Assistant e a exposição dos diagnósticos a um
editor por meio de um \emph{language server}, aproximando ainda mais a
verificação do momento da escrita.


\bibliographystyle{sbc}
\bibliography{sbc-template}

\end{document}
